\section{Discussion}\label{sec:discussion}

\subsection{What the evidence-grounded baseline establishes}

The principal result is not a new local LCZ map. It is the demonstration that a
locally assembled, interpretable evidence stack materially improves
spatially independent discrimination of global weak labels while the completed
audit exposes where those labels fail to represent local fabric. Morphology
alone carries substantial signal, but vegetation adds the largest early gain,
followed by green-blue and broader proxy context. This supports the literature's
move from class labels toward multidimensional urban form and vegetation structure
\citep{zhou2022UnderstaEffects,wang2023EffectivFactors,
bai2025ImpactBuilding,yu2026UrbanMorpholo}. At the same time, the importance of
WorldCover and Dynamic World variables warns that global land-cover semantics
may partly reproduce the target's own remote-sensing logic. The full-model
score cannot therefore be interpreted as purely independent local morphology.

The local 5 m surface model contributes surface roughness and relief at a much finer
resolution than the LCZ weak-label raster. Its scientific value lies in the
combination with exact footprints, floor-count proxies, and vegetation, not in
a claim that surface-model values are direct building heights. This distinction is
important because SVF and 3D urban-form studies demonstrate both the value and
the scale dependence of surrounding surface geometry
\citep{chen2010ViewFactor,middel2018ViewFactor,
wu2025ContrastUrban,yu2026UrbanMorpholo}. In Izmir, LCZ 3 combines high coverage
and vertical intensity, whereas LCZ 8 combines large-footprint form with the
highest median LST. These patterns are physically plausible, but only a manual
label audit can determine whether they represent locally coherent fabrics or
weak-label mixtures.

\subsection{Why 2SFCA is preferable to a green-space buffer}

The 2SFCA experiment changes the meaning of context. A Euclidean buffer answers
how much mapped green area lies within a circle. Network 2SFCA instead asks how
much supply is reachable through a road-cost graph after accounting for demand
competition. This is closer to an accessibility construct and avoids implying
that distance alone equals service. The positive primary-recipe increment for
LightGBM is modest and model-dependent, showing that the resulting field adds
some signal beyond local morphology and green-blue cover without becoming a
standalone accessibility measure.

Threshold results reinforce the need for sensitivity rather than a single
``correct'' catchment. LightGBM was marginally strongest at 1,200 m, whereas
the preregistered 800 m threshold remained positive across the primary model
comparison. Such
model dependence is expected when urban-form and temperature relationships
change with spatial support \citep{wu2022ScaleEffect,jeon2023ImpactsUrban,
li2026InvestigEffect}. We retain 800 m as the primary planning-scale proxy and
report the other thresholds rather than choosing the best after observing the
outcome. This avoids turning sensitivity analysis into hidden hyperparameter
optimization.

The municipal park comparison provides a useful but bounded validation. Strong
log-count and log-area correlations indicate that OSM supply follows the broad
district geography of the official maintenance inventory. Lower rank
correlations, sparse OSM names, and only 10.6\% exact-or-strong name candidates
show why this cannot validate individual polygons, public access, quality, or
completeness. The 2SFCA should therefore be interpreted as mapped green-supply
access, not a service-equity score.

\subsection{Spatial transfer is the central reviewer test}

Five spatial block folds produce stable average performance, but the LightGBM
LODO test reveals a 0.344 range in macro-F1. The hard transfers are not random noise; they
are evidence that district-specific combinations of coast, slope, building
stock, function, and label composition matter. The 0.115 difference between
mean spatial-fold and mean district-held-out macro-F1 is therefore more
informative than another decimal place in the headline score. It shows how a
model can interpolate local structure while struggling when an administrative
and morphological context is entirely absent from training.

This finding agrees with work showing that LCZ thermal meaning and urban-form
relationships vary across city, season, class, and analytical unit
\citep{kwok2019WellDoes,deng2024SensitivLocal,
wang2025SeasonalEffects,wu2025ContrastUrban}. A future external-city test would
be even harder. Before that step, the Izmir audit should oversample or at least
carefully inspect the districts with low LODO scores, while preserving the
original probability sample for unbiased reporting.

\subsection{Climate relevance and planning interpretation}

The LST contrast provides climate relevance without making LST a predictor.
The large Kruskal effect size and the warm LCZ 8 profile indicate that broad
weak-label fabric categories organize surface temperature. The correlation pattern is also
internally coherent: NDBI, building coverage, and road density are positive,
whereas NDVI and canopy volume are negative. Similar studies show that
morphology and vegetation effects can vary by LCZ, season, and scale
\citep{zhou2023IdentifyEffects,liu2025UnveilinDifferen,
wang2025SeasonalEffects}. Our bivariate results should not be converted into
causal prescriptions such as a fixed percentage of canopy producing a fixed
cooling benefit.

Planning translation should therefore be conditional. In large low-rise,
high-LST fabrics, priorities may include roof and yard heat mitigation,
continuous tree shade, and reducing exposed impervious service areas. In
compact low-rise fabrics, street-canyon shade, ventilation-sensitive planting,
and small-networked green spaces may be more relevant than large isolated
parks. In open low-rise and sparsely built areas, retaining vegetation and
preventing fragmentation may be more important than blanket densification.
These are mechanism-informed hypotheses for audited subtype panels, not current
policy prescriptions.

Nor does morphology equal vulnerability. LCZ and LST organize exposure, but
heat risk also depends on social sensitivity, mobility, occupation, housing,
and adaptive capacity \citep{ma2023InvestigUrban,mashhoodi2024UrbanForm,
iqbal2025DoesPerceive,turner2025DevelopmValidati}. The present model contains no
individual vulnerability data and should not be used to rank populations.
Future planning synthesis can join audited climate-morphology types to
area-level vulnerability only with separate governance, denominator, and
privacy checks.

\subsection{What a deep model must add}

A deep model is justified only if it extracts spatial configuration not already
captured by 68 tabular predictors. Figure-ground texture, street-building
interface, vegetation placement, and localized surface-model patterns are plausible
sources of additional information. Prior LCZ and urban-fabric work demonstrates
that CNN and multi-view representations can learn such structure
\citep{qiu2020MultilevFeature,yoo2020ImprovinLocal,
fang2024UrbanclaDeep}. But superiority must be established through identical
spatial folds, audited labels, branch ablations, probability calibration, and
attribution. A higher weak-label score without these checks would merely learn
the global label product more efficiently.

The current stopping rule is therefore a strength. It prevents a planned
architecture from being described as an implemented contribution. The completed
audit already shows that weak-label agreement is only 8.3\% in the primary
high-quality tier, so the key next test is not whether CLIMORFA beats Random
Forest by any weak-label margin, but whether raster and figure-ground branches
improve audited transfer, reduce systematic uncertainty, and reveal
interpretable subtype structure. SHAP and neural attributions must agree at the level of mechanism
families rather than being selected as decorative examples
\citep{xu2025CharacteThermal,han2025PerspectUndersta,
liu2026UnderstaNonlinea}.

\subsection{Limitations}

First, all classification results use global weak labels and four represented
built classes. The confidence threshold removes some uncertainty but can also
select easier, more homogeneous examples. Second, the manual audit is complete,
but the primary high-quality comparison contains only 96 four-class modeling
units and shows low agreement; it should not be generalized to a complete local
subtype census.
Third, the 5 m surface model is a surface model with uncertain vertical datum and no
terrain normalization, while floor-count height is a proxy. Fourth, Landsat LST
is a summer surface-temperature composite, not air temperature, thermal comfort,
or personal exposure. Fifth, missingness is partly structural, and median
imputation may blur the distinction between absence and unknown data.

Sixth, OSM functional context is not official zoning, and OSM park/garden supply
is neither a complete green inventory nor proof of access. The 2SFCA graph is
based on municipal road centrelines, not a fully audited pedestrian network;
explicit motorways were removed, but ramps or grade-separated segments may
remain under other codes. Seventh, endpoint-derived network measures are not a
fully planar topological graph, and street-frontage fields are road-buffer
proxies. Eighth, 500 m data are prepared but the complete 500 m model comparison
and optimized 125 m pipeline have not run. Finally, feature importance is
model-specific, correlations are bivariate, and no causal effect is estimated.

These limitations define a finite completion path: train raster and
figure-ground branches on identical audited spatial folds; run 125/250/500 m
analyses; calibrate uncertainty; perform branch-level XAI; and repeat the
climate and planning synthesis only for audited subtypes. The completed audit
is now an input to that next stage rather than an outstanding data gate.

